% Teaser (Figure 1): the close-loop arc. Self-contained TikZ, no external asset.
\begin{figure}[tbp]
  \centering
  \resizebox{\linewidth}{!}{%
  \begin{tikzpicture}[
      font=\small,
      box/.style={rounded corners=2pt, draw=black!55, line width=0.6pt,
                  align=center, inner sep=4pt, minimum height=15mm, text width=27mm},
      lab/.style={font=\scriptsize\itshape, text=black!55},
      arr/.style={-{Latex[length=2mm]}, line width=0.7pt, black!70},
    ]
    \node[box, fill=blue!6]   (p1) at (0,0)
      {\textbf{Silence gene $A$}\\[1pt] genome-wide\\ transcriptional\\ \emph{shockwave}};
    \node[box, fill=blue!6]   (p2) at (3.6,0)
      {\textbf{Reads $A$'s own}\\ \textbf{dependency}\\[1pt] transcriptome\\ $\to$ GeneEffect};
    \node[box, fill=orange!8] (p3) at (7.2,0)
      {\textbf{Names $A$'s SL}\\ \textbf{partners?}\\[1pt] CV2 anchored \cmark\\ CV3 both cold \xmark};
    \node[box, fill=green!8]  (p4) at (10.8,0)
      {\textbf{Generate the}\\ \textbf{shockwave}\\[1pt] AIVC/STATE\\ $+$ ESM2 identity};

    \draw[arr] (p1) -- (p2);
    \draw[arr] (p2) -- (p3);
    \draw[arr] (p3) -- (p4);
    % closing arrow: predicted shockwave feeds the partner question for unseen genes
    \draw[arr, dashed, black!60] (p4.south) .. controls (10.8,-1.5) and (7.2,-1.5) ..
      (p3.south) node[midway, below, lab, text width=42mm]
      {a \emph{predicted} shockwave should carry the same\\ partner signal --- now for genes never screened};

    \node[lab, below=1mm of p1] {exp01 foundation};
    \node[lab, below=1mm of p2] {exp01--03};
    \node[lab] at ([yshift=-1mm]p3.south) {};
  \end{tikzpicture}}
  \caption{The transcriptional shockwave carries the partner. Perturbing a gene
  produces a genome-wide response that predicts the gene's own dependency
  (exp01). The same observed response also ranks a gene's synthetic-lethal
  partners when one partner is anchored by a real profile (CV2), but not when
  both genes are unscreened (CV3). Our contribution is a virtual-cell framework
  that \emph{generates} this response from a protein-language-model gene
  identity, so the partner signal can in principle reach pairs neither gene was
  screened for (dashed loop; results preliminary).}
  \label{fig:teaser}
\end{figure}
